% !TEX root = ../cubic-moduli-descent.tex
\section{Hesse 中心の正則化と一因子周期の全係数}
\label{mbhp:sec:hesse-phase}

本節では、補題\ref{mbpd:lem:hodge-line}の比較直線と、
定理\ref{mbpl:thm:global-line-comparison}の自然な関係直線同期を用いる。
その指定周期の三乗写像を、Jacobian 上の Weierstrass 奇函数によって計算する。
中心元の既知の Hesse 式を使うが、中心直線の flex における正則化、
比較の原点での規格化、および周期 torsor の全係数の決定が本節の内容である。
以下の境界への延長は計算のための一因子模型であり、
全 $\OO$ 認識の領域を境界へ広げる主張ではない。

\subsection{固定した Hesse 曲線上の中心直線}

$K\supset\C$、$\mu^3\ne1$ とし、
\[
 E=V(F_\mu),\qquad
 F_\mu=x^3+y^3+z^3-3\mu xyz,\qquad
 o=(1:-1:0),\qquad L=\OO_E(1)=\OO_E(3o)
\]
を固定する。平行移動点のパラメータ曲線を $P=E$ と書く。
$p=(a:b:c)\in P$ に対し、族 $\mathcal S$ の関係は
\[
 cx^2+ayz+bzy,\qquad cy^2+azx+bxz,\qquad cz^2+axy+byx
\]
である。関係の共通 scalar を除いたこの表示は $P$ 上の関係部分束を定める。
以下、$x,y,z$ は代数の生成元、$a,b,c$ はパラメータである。
$C=x^3+y^3+z^3$、$U=xyz+yzx+zxy$、
$V_0=xzy+zyx+yxz$ と書く。

\begin{lemma}[九つの flex を含む中心直線の正則化]
\label{mbhp:lem:regularized-center}
中心式\eqref{mbpl:eq:hesse-center}について、
$abc\ne0$ 上の
\begin{equation}\label{mbhp:eq:regularized-center}
 \widetilde\Omega=\frac{\Omega_H}{abc}
\end{equation}
は、$P$ 全体で非零な切断
$\mathcal S_3\otimes\OO_P(1)$ に延長する。
その像は直線部分束
\[
 \mathcal C\simeq\OO_P(-1)=\OO_P(-3o)\longrightarrow\mathcal S_3
\]
であり、全幾何ファイバーで中心正則三次を与える。
その商は $(E,L,t_p)$ の相対 THCR である。
点方程式の固定直線を
$\mathcal J=KF_\mu\otimes\OO_P$ とすると、基本開の自然な比較は
束の射
\begin{equation}\label{mbhp:eq:extended-theta}
 \vartheta:\mathcal C\longrightarrow\mathcal J,\qquad
 \vartheta(\widetilde\Omega)=3(a-b)F_\mu
\end{equation}
へ延長する。その零因子は
$E[2]\setminus\{o\}$ の三つの点の和である。
\end{lemma}
\begin{proof}
各ファイバーは、非 flex では通常の滑らかな楕円 Sklyanin 代数である。
通常の AS-regular 性の例外十二点は、三つの座標頂点と
$a^3=b^3=c^3$ の九点であり、$\mu^3\ne1$ の $E_\mu$ 上にはない。
flex では $p\in E[3]$ なので $t_p^*L\simeq L$ であり、
$t_p$ は $H^0(E,L)$ 上の線型持上げを持つ。
この持上げによる $\operatorname{Sym}H^0(E,L)$ の Zhang twist を取る。
平面三次の次数二イデアルは零なので、その二次関係は
$(E,L,t_p)$ の幾何二次関係そのものである。
従って $a=0$、$b=0$、$c=0$ の全九 flex で、
$\mathcal S_p$ はこの多項式環の Zhang twist に一致する。
特に全ファイバーは三次元 quadratic AS-regular domain であり、
Hilbert 級数は $(1-t)^{-3}$ である
\cite[Theorems~6.6, 6.8]{GRSklyaninATV}
\cite[Section~1]{GRSklyaninST}。
有限関係行列の階数は各次数で一定なので、
滑らかな被約曲線 $P$ 上の $\mathcal S_n$ は局所自由である。
この議論は中心商の幾何を先に仮定していない。

次数三の不変関係は
\[
 cC+aU+bV_0=0
\]
である。$b\ne0$ で $V_0$ を消去し、
$a^3+b^3+c^3=3\mu abc$ を用いると
\begin{equation}\label{mbhp:eq:division-certificate}
 \Omega_H
   =3ac\bigl((a^2-\mu bc)C+(\mu ab-c^2)U\bigr)
       \quad\text{in }\mathcal S_3
\end{equation}
を得る。従って
\begin{equation}\label{mbhp:eq:regular-chart}
 \widetilde\Omega
   =\frac3b\bigl((a^2-\mu bc)C+(\mu ab-c^2)U\bigr).
\end{equation}
$a\ne0$ では $U$ を消去して
\[
 \widetilde\Omega
   =-\frac3a\bigl((b^2-\mu ac)C+(\mu ab-c^2)V_0\bigr)
\]
を得る。$\{a\ne0\}$ と $\{b\ne0\}$ は $P$ を覆うので、
この二式は九つの flex の全てを含む。
これらは稠密開で一致するため、局所自由な $\mathcal S_3$ 上で貼り合わさる。
係数は $(a,b,c)$ に関して斉次次数一であり、
正確には $\mathcal S_3\otimes\OO_P(1)$ の切断を得ている。
ここでの除算は $\mathcal S_3$ 内の恒等式である。
パラメータ torsion を持ち得る最大可換商の次数三で
$abc$ を消去しているのではない。

$c=0$、$b=-\zeta a$、$\zeta^3=1$ では
\[
 \widetilde\Omega=-3a\zeta^{-1}C+3\mu aU.
\]
補題\ref{mbpl:lem:hesse-center}の不変部分の計算はこの点でも適用でき、
唯一の不変関係は $aU+bV_0$ である。
よって $C,U$ はその商で独立であり、上の元は非零である。
非 flex では同補題の非零性の証明を用いる。
$a=0$ では
$\widetilde\Omega=-3\mu cC-3c^2U/b$ であり、
唯一の不変関係は $cC+bV_0$ なので非零である。
$b=0$ では
$\widetilde\Omega=3\mu cC+3c^2V_0/a$、
唯一の不変関係は $cC+aU$ であり、やはり非零である。
既知の中心式
\cite[Section~2, p.~3]{GRSklyaninDLL}
\cite[p.~35, Remark~(1)]{GRSklyaninST}
と $\mathcal S_4$ の局所自由性により、
中心性は稠密開から全体に延長する。
各ファイバーが domain であるため、非零中心元は正則である。

幾何的関係の評価は $\mathcal S$ から相対 THCR への射を与える。
THCR の各次数は局所自由であり、$\widetilde\Omega$ の像は稠密開で零なので
全体で零である。次数三の核は階数一で、
上で得た非零中心元がこれを生成する。
中心正則性による商の Hilbert 級数と THCR の Hilbert 級数は一致する。
評価の次数一生成性から全射であり、従って全次数で同型となる。

$F_H=abcF_\mu$ なので、
式\eqref{mbpl:eq:hesse-comparison}は
式\eqref{mbhp:eq:extended-theta}を与える。
この射は $\OO_P(-1)\to\OO_P$ という斉次性とも整合する。
$o$ を原点とする反転は $(a:b:c)\mapsto(b:a:c)$ である。
その固定点は $o$ と直線 $a=b$ の三交点であり、
標数零の滑らかな楕円曲線の二捩れ点は被約である。
従って主張した零因子を得る。
\end{proof}

この族への相対 window の適用には
系\paperref{main}{mbrn:cor:factor-boundary}を用いる。
ここで $\mathcal C$ は指定した THCR 商の核の直線であり、
中心の次数三全体の次元が境界でも一であるとは主張しない。
特に最上 Koszul syzygy 直線は実際の
$(W\otimes R_2)\cap(R_2\otimes W)$ のまま保持する。
二捩れ点で退化することのある $\det W$ への交代化は使わない。
中心商、有限 Koszul 模型、指定周期はこの補題の全パラメータ上で存在する。

\subsection{比較直線の原点を含む延長}

以下の符号規約を固定する。
対角の余法束の同型は
\begin{equation}\label{mbhp:eq:conormal-sign}
 I_\Delta/I_\Delta^2\longrightarrow\Omega^1_E,\qquad
 \overline{\operatorname{pr}_2^*f-\operatorname{pr}_1^*f}\longmapsto df
\end{equation}
とする。これが補題\ref{mbpd:lem:hodge-line}の同型の符号を定める。

\begin{lemma}[比較直線と同期前の三乗写像]
\label{mbhp:lem:global-comparison-line}
$E_x\times P$ 上の線束
$M=\OO(\{x=p\}-\{x=o\})$ を用い、$F_M=-\otimes M$ とする。
$R_p=T_{\OO_E}(t_p)_*(-\otimes L)$、
$R_0=T_{\OO_E}(-\otimes L)$ に対する比較射の直線 $\mathcal V$ は、
$p=o$ を含めて
\begin{equation}\label{mbhp:eq:global-lines}
 \mathcal V\simeq\lambda\otimes\OO_P(o),\qquad
 \widetilde q:\mathcal V^{\otimes3}
       \xrightarrow{\sim}\mathcal J\otimes\mathcal C^\vee,
 \qquad \lambda=H^0(E,\Omega_E^1)
\end{equation}
を満たす。
$p\ne o$ 上の $\mathcal V\simeq\lambda$ は
補題\ref{mbpd:lem:hodge-line}のものと一致する。
$\vartheta$ で同期し、標準切断による
$\lambda\to\lambda\otimes\OO_P(o)$ を三乗して得る
有理写像 $q:\lambda^{\otimes3}\dashrightarrow\OO_P$ の因子は
\begin{equation}\label{mbhp:eq:period-divisor}
 \operatorname{div}(q)=3[o]
       -\sum_{\tau\in E[2]\setminus\{o\}}[\tau].
\end{equation}
\end{lemma}
\begin{proof}
$E_x\times E_y\times P$ 上で $\Delta=(y=x)$、
$\Gamma=(y=x+p)$ と置く。
式\eqref{mbpd:eq:kernel-ratio}と同じ計算により、核の比は
\[
 N=\OO(\Gamma-\Delta)\otimes
       \operatorname{pr}_{x,p}^*M\otimes
       \operatorname{pr}_{y,p}^*M^{-1}
\]
である。$p\ne o$ では相対 Picard intertwiner の計算を使い、
$p=o$ では二つの核と $M_o=\OO_E$ が一致する。
従って全パラメータで $N$ のファイバーは自明な線束である。
proper cohomology と底変更、または同じ相対 intertwiner の構成により、
$g:E\times E\times P\to P$ に対して
$\mathcal V=g_*N$ は直線束、$g^*\mathcal V\to N$ は同型となる。
ここでは滑らかな被約のパラメータ曲線上で証明しており、
非被約底上の結論を幾何ファイバーの一致だけから出していない。

対角上では二つの $M$ は標準的に相殺する。
$\Gamma$ と $\Delta$ の交わりは正確に $p=o$ なので、
\[
 N|_{\Delta\times P}
     \simeq\Omega_E^1\boxtimes\OO_P(o).
\]
これが $\mathcal V$ の主張を与える。
$p\ne o$ では $\OO_P(o)$ の標準切断で自明化するため、
既存の Hodge 同型とも一致する。

中心商の相対 window は、まだ同期していない二つの指定周期
\[
 R_p^3\simeq[2]\otimes\mathcal C,\qquad
 R_0^3\simeq[2]\otimes\mathcal J
\]
を与える。比較射を三回合成してこれらを適用すると、
$\mathcal V^3\to\mathcal J\otimes\mathcal C^\vee$ を得る。
各ファイバーで比較射と周期はいずれも可逆なので、この写像は同型である。
これは系\paperref{main}{mbrn:cor:factor-boundary}の実際の window 周期であり、
途中で scalar を取り替えたものではない。
$\vartheta$ は $\mathcal J\otimes\mathcal C^\vee$ の切断であるから、
同期後の $q$ は $\widetilde q$ をこの切断で割った写像になる。
$\lambda\to\mathcal V$ の零因子は $[o]$、
$\vartheta$ の零因子は前補題の三つの二捩れ点である。
これから式\eqref{mbhp:eq:period-divisor}を得る。
\end{proof}

\subsection{Weierstrass 奇函数による全係数}

楕円曲線の原点付き Weierstrass 表示と微分を
\[
 y^2+a_1xy+a_3y=x^3+a_2x^2+a_4x+a_6,\qquad
 v_\omega=2y+a_1x+a_3,\qquad \omega=dx/v_\omega
\]
と書く。この表示とその座標変換則には
\cite[Chapter~III, \S1, Table~3.1; \S5]{NFMSilvermanArithmetic}
を用いる。$v_\omega$ は反転で符号を変え、因子は
\[
 \operatorname{div}(v_\omega)
   =\sum_{\tau\in E[2]\setminus\{o\}}[\tau]-3[o]
\]
である。

\begin{theorem}[自然な同期に対する一因子周期の全係数]
\label{mbhp:thm:period-formula}
式\eqref{mbhp:eq:conormal-sign}の規約を用いる。
定理\ref{mbpl:thm:global-line-comparison}の自然な同期に対し、
補題\ref{mbpd:lem:hodge-line}の三乗写像は
\begin{equation}\label{mbhp:eq:period-formula}
 q(\omega^{\otimes3})=-\frac1{v_\omega(p)}.
\end{equation}
これは $p,2p,3p\ne0$ の全基本楕円開上の式であり、
普遍族からの任意の許された底変更に両立する。
より内在的には、
$\mathfrak v(p)=v_\omega(p)\omega^{\otimes3}$
という $\lambda^{\otimes3}$ の可逆切断について
$q(\mathfrak v(p))=-1$ である。
\end{theorem}
\begin{proof}
まず固定した $E_\mu$ 上で証明する。
前補題の因子公式により、
$q(\omega^{\otimes3})v_\omega$ は $p$ に依存しない非零定数である。
その値を原点で求める。
$a=1,b=-1,c=0$ では代数は可換多項式環であり、
式\eqref{mbhp:eq:regular-chart}から
\begin{equation}\label{mbhp:eq:origin-curvatures}
 \widetilde\Omega_o=-3C+3\mu U=-3F_\mu,\qquad
 \vartheta(\widetilde\Omega)_o=6F_\mu
\end{equation}
となる。ここでは可換代数の中で $U=3xyz$ である。
曲率を $f$ から $cf$ へ変えると指定周期は
$\pi_{cf}=c^{-1}\pi_f$ となる
（命題\ref{prop:nfm-period-equation-frame-change}）。
原点では両側の商代数と window の section 切断が同一であり、
$M_o=\OO_E$ の標準自明化により比較射は恒等射 $\eta_o$ である。
従って、同期済みだが $\mathcal V_o$ 上で評価した三乗写像は
\begin{equation}\label{mbhp:eq:origin-cube}
 q_{\mathcal V,o}(\eta_o^{\otimes3})
        =\frac{1/6}{-1/3}=-\frac12.
\end{equation}
同一の quotient/window に二つの方程式 frame を入れた比較なので、
ここに未指定の window scalar はない。

Weierstrass の局所径数を $t=-x/y$ とすれば
\[
 \omega=dt+O(t)\,dt,\qquad
 v_\omega=-2t^{-3}+O(t^{-2}).
\]
$\lambda\to\mathcal V$ の一次項も必要である。
曲線上の原点とは異なる点の局所座標を $u$、出力座標を $z$ とし、
平行移動を
$z=u+A(u)t+O(t^2)$、$\omega=du/A(u)$ と書く。
この近傍では $M$ の標準有理切断を用い、
核の比線束の局所 frame を
\[
 e=\frac{z-u}{z-u-A(u)t+O(t^2)}
\]
と取れる。$t=0$ でこれは恒等射 $\eta_o$ に一致する。
対角上で分子の余法類を式\eqref{mbhp:eq:conormal-sign}により $du$ へ送り、
分母を $-A(u)t+O(t^2)$ へ制限すると、
$\mathcal V\to\lambda$ の有理同型は
$e\mapsto-\omega/t+O(1)$ となる。
従って逆向きの正則射は
\[
 \omega\longmapsto-t\eta_o+O(t^2).
\]
式\eqref{mbhp:eq:origin-cube}と合わせて
$q(\omega^3)=t^3/2+O(t^4)$ を得る。
$v_\omega$ の主部を掛ければ定数は $-1$ となり、
式\eqref{mbhp:eq:period-formula}が従う。

この議論をまず $K=\C(\mu)$ に適用する。
滑らかな普遍 Hesse 被覆上では両辺は正則な線束の射であり、
稠密な一般点での一致は全基本開上の一致を与える。
中心線比較と window は表示の同型に自然である。
また Weierstrass 座標変更
$x=u^2x'+r$、$y=u^3y'+su^2x'+t_0$ に対して
$v_\omega=u^3v_{\omega'}$、$\omega=u^{-1}\omega'$ なので、
$v_\omega\omega^3$ は不変である。
従って普遍 Hesse 被覆の重なりでも式は一致し、
補題\ref{mbpl:lem:hesse-cover}の fpqc 降下により
全基本楕円族の式となる。
最後にこの普遍族上の射の等式を引き戻す。
これにより非被約な底変更にも対応し、
任意底上で幾何点だけを比較する議論を避けている。
\end{proof}

\begin{corollary}[全基本開上の Kummer torsor]
\label{mbhp:cor:global-kummer}
原点を選ばない genus-one 曲線族についても、
Jacobian $J=\Pic^0(E/S)$ とその点 $p$ を用いて
$\mathfrak v(p)\in\Gamma(S,\lambda^{\otimes3})^\times$ が内在的に定まる。
自然な同期に対する因子周期 torsor は
\begin{equation}\label{mbhp:eq:cube-torsor}
 T\simeq\{\xi\in\lambda^\times:
              \xi^{\otimes3}=\mathfrak v(p)\}.
\end{equation}
特に Hodge frame $\omega$ を選んだ底上で
\[
 \gamma=[v_\omega(p)]\in H^1(S,\mu_3).
\]
この記述は因子の同型、共通 grading scalar、
中心生成元とその自然な像の同時変更、および任意底変更に不変である。
\end{corollary}
\begin{proof}
$E$ と $J$ の Hodge 線は標準的に一致する。
Weierstrass 変換則により $\mathfrak v$ は Jacobian 上で定義され、
原点と非零二捩れ点を除けば可逆である。
定理の式より、周期適合条件 $q(\eta^3)=1$ は
$\eta^3=-\mathfrak v(p)$ に等しい。
$\xi=-\eta$ とすれば式\eqref{mbhp:eq:cube-torsor}を得る。
これは $\mu_3$-同変である。
Hodge frame の変更は函数を三乗倍するだけであり、
中心方程式の同時変更は二つの周期を同じ比で変更する。
その他の自然性は普遍族上の構成と前定理の降下による。
\end{proof}

\begin{corollary}[Hesse 座標での有理式]
\label{mbhp:cor:hesse-expression}
$a=1$ の chart 上で
$\omega_H=dc/(3b^2-3\mu c)$ と置く。
これは $E_\mu$ の不変微分であり、
$l=a+b+\mu c$ と書けば
\begin{equation}\label{mbhp:eq:hesse-phase}
 \begin{split}
 v_{\omega_H}(p)&=
       9(1-\mu^3)\frac{a-b}{a+b+\mu c},\\
 q(\omega_H^{\otimes3})&=
       \frac{a+b+\mu c}{9(\mu^3-1)(a-b)},\\
 \gamma&=\left[
       9(1-\mu^3)\frac{a-b}{a+b+\mu c}\right].
 \end{split}
\end{equation}
\end{corollary}
\begin{proof}
直線 $l=0$ は $o$ における接線であり、$E$ 上の因子は $3[o]$ である。
したがって $(a-b)/l$ は $v_{\omega_H}$ と同じ因子を持つ。
原点近傍で Hesse 方程式を解くと
\[
 b=-1-\mu c+\frac{\mu^3-1}{3}c^3+O(c^4),\qquad
 l=\frac{\mu^3-1}{3}c^3+O(c^4),\qquad
 \omega_H=(1/3+O(c))\,dc.
\]
$t=c/3+O(c^2)$ に対する主部
$v_{\omega_H}=-54c^{-3}+O(c^{-2})$ から係数が決まる。

別の直接検算として
\[
 X=\frac{3(\mu^3-1)c}{l},\qquad
 V=\frac{9(1-\mu^3)(a-b)}{l}
\]
を用いれば、Hesse 方程式の剰余として
\[
 \frac{dX}{\omega_H}=V,\qquad
 V^2=4X^3-27\mu^2X^2+
       54\mu(\mu^3-1)X-27(\mu^3-1)^2
\]
が成立する。従って $Y=V/2$ は
$\omega_H=dX/(2Y)$ の Weierstrass 表示を与える。
残りの式は前定理と前系による。
\end{proof}

\begin{corollary}[既存の単元表示と境界剰余]
\label{mbhp:cor:unit-and-residue}
$j\ne0,1728$ 上で $h=c_6/c_4\in\lambda^2$ と置くと、
系\ref{mbpd:cor:unit-phase}の単元と frame は
\[
 u=\frac{h^3}{\mathfrak v(p)^2},\qquad
 e=-\frac{h^2}{\mathfrak v(p)}
\]
である。Hodge frame に関する $h=H\omega^2$、
$\mathfrak v=V\omega^3$ では
$[u]=[H^3/V^2]=[V]=\gamma$ となる。
また良い還元を持つ一径数族で $p$ が非零二捩れ因子と横断的に交われば、
自然な位相の Kummer 剰余は $+1\in\mathbb Z/3$ である。
\end{corollary}
\begin{proof}
$q(\mathfrak v)=-1$ を
$e=(q\otimes\mathrm{id})(h^2)$、$u=q^{\otimes2}(h^3)$ へ代入する。
剰余の主張は $v_\omega$ が非零二捩れ点で単純零を持つことによる。
従って命題\ref{mbpd:prop:boundary-phase}の非零剰余と一致する。
\end{proof}

本節が決定したのは、自然な関係直線同期に対する一因子の
Picard 補正周期である。
これを積に適用する際の scalar gerbe、対角位相商、
および選んだ比較核型の条件は定理\ref{mbpd:thm:coherent-corrections}に従う。
この有理式だけから全核型を同定したり、
表示 gerbe や比較核の大域実現障害が消えると結論したりはしない。

中心の除算恒等式、Hesse--Weierstrass 変換、原点での形式展開、および
関連する period multiplier の例の有限多項式計算は
\texttt{computations/verify\_multiblock\_hesse\_phase.py} で再検算できる。
この計算は係数と恒等式を確認するものであり、
window による圏比較、Hom 直線、または fpqc 降下の機械的証明ではない。
